% How much of the condition survives the frozen VAE. Plotted from vae.dat,
% written by scripts out of results/. Same drafting as the other plots.
\begin{figure}[t]
\centering
\begin{tikzpicture}
\begin{axis}[draft, height=42mm,
  xlabel={cutoff $r$}, ylabel={correlation after round trip},
  xmin=-0.03, xmax=1.03, ymin=0, ymax=1.05,
  legend pos=north east]
  \addplot[mark=square*, black] table[x=r, y=native] {vae.dat};
  \addlegendentry{condition built at 512}
  \addplot[mark=o, black, densely dashed] table[x=r, y=up256] {vae.dat};
  \addlegendentry{built at 256, upsampled}
\end{axis}
\end{tikzpicture}
\caption{Correlation between the condition fed to the autoencoder and the same
condition re-extracted after an encode--decode round trip, the first 24
FFHQ-512 images in file order, which are not the test split. Above
$r{\approx}0.5$ the condition is effectively gone. Building the condition at 256
and upsampling puts the
structure at lower absolute frequency and preserves far more, which is why
comparisons against older pipelines need the source resolution stated.}
\label{fig:vae}
\end{figure}
